% PDF page 372
\source{gilbarg-trudinger:page-0372}

(i) If $Q$ is elliptic with principal coefficients $a^{ij}$ depending only on $p$, then clearly we can take $a_*^{ij}=a^{ij}$ and $c_i=0$;

(ii) The Euler--Lagrange equation corresponding to a multiple integral of the form

\begin{equation}
\int_{\Omega} F(x,u,|Du|)\,dx,
\tag{15.3}
\end{equation}

where $F\in C^2(\Omega\times\mathbb{R}\times\mathbb{R})$, $D_tF\neq 0$, $t=|p|$, can be written as

\begin{equation}
\Delta u+\{(|Du|D_{tt}F/D_tF)-1\}|Du|^{-2}D_i u D_j u D_{ij}u
+ (|Du|^2D_{tz}F + D_i u D_{tx_i}F - |Du|D_zF)/D_tF=0,
\tag{15.4}
\end{equation}

so that the decomposition (15.2) is valid with

\[
a_*^{ij}=\delta^{ij},\qquad c_i=\{(|p|D_{tt}F/D_tF)-1\}\,|p|^{-2}p_i
\]

provided $c_i\in C^1(\Omega\times\mathbb{R}\times\mathbb{R}^n)$, $i=1,\ldots,n$.

(iii) In the special case of two variables we can write

\begin{equation}
 a^{ij}=(\Tr-\Estar)\delta^{ij}+\tfrac12(p_i d_j+d_i p_j)
\tag{15.5}
\end{equation}

where

\[
\Tr=\operatorname{trace}[a^{ij}]=a^{11}+a^{22}
\]
\[
\Estar=\mathcal{E}/|p|^2=a^{ij}p_ip_j/|p|^2,
\]
\[
d_1=[(a^{11}-a^{22})p_1+2a^{12}p_2]/|p|^2,
\]
\[
d_2=[2a^{12}p_1+(a^{22}-a^{11})p_2]/|p|^2.
\]

Hence for $n=2$, the equation (15.1) is equivalent to the equation

\begin{equation}
\Delta u+c_iD_j uD_{ij}u+b^*=0
\tag{15.6}
\end{equation}

where $c_i=d_i/(\Tr-\Estar)$, $b^*=b/(\Tr-\Estar)$; the decomposition (15.2) is clearly valid for equation (15.6) with $a_*^{ij}=\delta^{ij}$.

The basic idea in our treatment of gradient bounds, which goes back as far as Bernstein's work [BE 1], involves differentiation of equation (15.1) with respect to $x_k$, $k=1,\ldots,n$, followed by multiplication by $D_k u$ and summation over $k$. The maximum principle is then applied to the resulting equation in the function $v=|Du|^2$. By (15.2) we can write equation (15.1) in the form

\[
a_*^{ij}(Du)D_{ij}u+\tfrac12 c_i(x,u,Du)D_iv+b(x,u,Du)=0.
\]

% PDF page 373
\source{gilbarg-trudinger:page-0373}



\section*{15.1. A Maximum Principle for the Gradient}

Assuming that the solution $u \in C^3(\Omega)$ we then obtain, by differentiation with respect to $x_k$,

\begin{equation*}
a_*^{ij}D_{ijk}u + D_{p_i}a_*^{ij}D_{ik}uD_{ij}u+\frac12(c_iD_{ik}v+D_kc_iD_iv)
+ D_{p_i}bD_ku + D_kuD_zb + D_{x_k}b = 0.
\end{equation*}

Multiplying by $D_ku$ and summing over $k$, we thus have

\begin{equation}
(15.7)\qquad -2a_*^{ij}D_{ik}uD_{jk}u+a_*^{ij}D_{ij}v+\bigl(D_{p_i}a_*^{ij}D_{ij}u+D_{p_i}b+D_kuD_kc_i\bigr)D_iv+2\delta bv=0
\end{equation}

where $\delta$ is the differential operator, acting on $C^1(\Omega\times \mathbb{R}\times \mathbb{R}^n)$, defined by

\begin{equation}
(15.8)\qquad \delta=D_z+|p|^{-2}p_iD_{x_i}.
\end{equation}

We next apply the following consequence of Schwarz's inequality

\begin{equation}
(15.9)\qquad a_*^{ij}D_{ik}uD_{jk}u\ge (a_*^{ij}D_{ij}u)^2/\operatorname{trace}[a_*^{ij}]
=(b+\tfrac12c_iD_iv)^2/\operatorname{trace}[a_*^{ij}].
\end{equation}

and hence obtain

\begin{equation}
(15.10)\qquad a^{ij}D_{ij}v+B_iD_iv\ge 2\left(\frac{b^2}{\mathcal{J}_*}-v\,\delta b\right)
\end{equation}

where $\mathcal{J}_*=\operatorname{trace}[a_*^{ij}]$ and

\[
B_i=D_{p_i}a_*^{ij}D_{ij}u+D_{p_i}b+D_kuD_kc_i-2\mathcal{J}_*^{-1}bc_i-\frac12\mathcal{J}_*^{-1}c_ic_jD_jv.
\]

Consequently, if the inequality

\begin{equation}
(15.11)\qquad |p|^2\delta b\le \frac{b^2}{\mathcal{J}_*}
\end{equation}

holds in $\Omega\times \mathbb{R}\times \mathbb{R}^n$, we obtain immediately from the classical maximum principle, Theorem 3.1, that $\sup_\Omega v=\sup_{\partial\Omega}v$. In order to extend this result to $C^2(\Omega)$ solutions we write equation (15.7) in the divergence form

\begin{equation*}
-2a_*^{ij}D_{ik}uD_{jk}u+D_i(a^{ij}D_jv)+\bigl(D_{p_i}a_*^{ij}D_{ij}u-D_ja^{ij}+D_kuD_kc_i\bigr)D_iv
+2D_k(bD_ku)-2\Delta ub=0,
\end{equation*}

which has the corresponding integral form

\begin{equation*}
2\int_\Omega \eta a_*^{ij}D_{ik}uD_{jk}u\,dx+\int_\Omega (a^{ij}D_jv+2bD_iu)D_i\eta\,dx
\end{equation*}

\begin{equation}
(15.12)\qquad -\int_\Omega \bigl\{(D_{p_i}a_*^{ij}D_{ij}u-D_ja^{ij}+D_kuD_kc_i)D_iv\bigr\}\eta\,dx
+\int_\Omega 2b\eta\,\Delta u\,dx=0
\end{equation}


% PDF page 374
\source{gilbarg-trudinger:page-0374}

for all $\eta \in C_0^1(\Omega)$ (see equation 13.17). An approximation argument as in Section
13.3 shows that equation (15.12) continues to be valid for $u \in C^2(\Omega)$. Integrating
the term
\[
\int_\Omega 2b\eta\,\Delta u\,dx
\]
by parts, and then proceeding as above for the case of $C^3(\Omega)$
solutions we obtain the weak inequality

\begin{equation}
\tag{15.13}
D_i(a^{ij}D_j v)+(B_i-D_ja^{ij})D_iv \ge 2(b^2/\mathcal{J}_* - v\,\delta b)
\end{equation}

for $v \in C^1(\bar{\Omega})$ and the estimate $\sup_\Omega v = \sup_{\partial\Omega} v$ follows from the maximum principle,
Theorem 8.1.

We have therefore proved the following gradient maximum principle.

\medskip

\noindent\textbf{Theorem 15.1.} \textit{Let $u \in C^2(\Omega) \cap C^1(\bar{\Omega})$ satisfy equation (15.1) in the bounded
domain $\Omega$ and suppose that $Q$ is elliptic in $\Omega$ with coefficients satisfying (15.2) and
(15.11). Then we have the estimate}

\begin{equation}
\tag{15.14}
\sup_{\Omega} |Du| = \sup_{\partial\Omega} |Du|
\end{equation}

\noindent It is clear from the above proof that conditions (15.2) and (15.11) in the
hypotheses of Theorem 15.1 need only hold for $|z| \leq \sup_\Omega |u|$. Furthermore, if they
also hold for $|p| \geq L$ for some constant $L$, we obtain in place of (15.14) the estimate

\begin{equation}
\tag{15.15}
\sup_{\Omega} |Du| \leq \max \left\{ \sup_{\partial\Omega} |Du|,\, L \right\}.
\end{equation}

\noindent The estimate (15.15) follows by applying Theorem 15.1 in the domain $\Omega_L =
\{x \in \Omega \mid |Du| > L\}$.

\noindent Note that when Theorem 15.1 is applied to examples (ii) and (iii) above where
$a_*^{ij} = \delta^{ij}$ the condition (15.11) reduces to

\begin{equation}
\tag{15.16}
|p|^2\,\delta b \leq \frac{b^2}{n}
\end{equation}

\noindent In particular if $b=0$, we obtain a gradient maximum principle.

\bigskip

\noindent\textbf{15.2. \ The General Case}

\medskip

The technique to be employed in this section is basically a modification of that in
the preceding section through the use of an auxiliary function. Assuming initially
that $u$ is a $C^2(\Omega)$ solution of equation (15.1), we set $m=\inf_\Omega u$, $M=\sup_\Omega u$ and let $\psi$
be a strictly increasing function in $C^3[\bar m, \bar M]$ with $m=\psi(\bar m)$, $M=\psi(\bar M)$. Defining

% PDF page 375
\source{gilbarg-trudinger:page-0375}



the function $\bar u$ by $u=\psi(\bar u)$, so that

\[
D_i u=\psi' D_i \bar u
\]
\[
D_{ij}u=\psi'' D_i \bar u D_j \bar u+\psi' D_{ij}\bar u,
\]

we have the equation (see also (14.4)).

\begin{equation}
\psi' a^{ij}(x,u,Du)D_{ij}\bar u+b(x,u,Du)+\frac{\psi''}{(\psi')^2}\mathcal E(x,u,Du)=0.
\tag{15.17}
\end{equation}

Let us now write $v=|Du|^2$, $\bar v=|D\bar u|^2$ and apply the operator $D_k u D_k$ to the equation
(15.17). We obtain thus

\begin{align}
&-a^{ij}D_{ik}\bar u D_{jk}\bar u+\frac12 a^{ij}D_{ij}\bar v+\frac{1}{2(\psi')^2}\Bigl(\psi' D_{p_l}a^{ij}D_{ij}\bar u+D_{p_l}b
\notag\\
&\qquad+\frac{\psi''}{(\psi')^2}D_{p_l}\mathcal E\Bigr)D_l v+\Bigl(\psi' \delta a^{ij}D_{ij}\bar u+\delta b+\frac{\psi''}{(\psi')^2}\delta \mathcal E\Bigr)\bar v
\tag{15.18}\\
&\qquad-\frac{\psi''}{(\psi')^2}\Bigl(b+\frac{\psi''}{(\psi')^2}\mathcal E\Bigr)\bar v-\frac{1}{\psi'}\Bigl[\frac{\psi''}{(\psi')^2}\Bigr]'\delta \bar v=0.
\notag
\end{align}

Next, letting $\delta$ be the operator acting on $C^1(\Omega\times \mathbb R\times \mathbb R^n)$ defined by

\begin{equation}
\delta=p_iD_{p_i}
\tag{15.19}
\end{equation}

and introducing the function

\begin{equation}
\omega=\frac{\psi''}{(\psi')^2}\in C^1[\bar m,\bar M],
\tag{15.20}
\end{equation}

we obtain, with the aid of the relation

\[
D_l v=2\psi''\bar v D_l\bar u+(\psi')^2D_l\bar v,
\]

the following equation:

\begin{align}
&-a^{ij}D_{ik}\bar u D_{jk}\bar u+\frac12 a^{ij}D_{ij}\bar v+\frac12\Bigl(\psi' D_{p_i}a^{ik}D_{jk}\bar u+D_{p_i}b+\omega D_{p_i}\mathcal E\Bigr)D_i \bar v
\notag\\
&\qquad+\psi'\bar v(\omega \delta a^{ij}+\delta a^{ij})D_{ij}\bar u
\tag{15.21}\\
&\qquad+\Bigl\{\frac{\omega'}{\psi'}\mathcal E+\omega^2(\delta-1)\mathcal E+\omega(\delta \mathcal E+(\delta-1)b)+\delta b\Bigr\}\bar v=0.
\notag
\end{align}

Equation (15.21) can be further generalized by combining it with equation (15.17).
Letting $r$ and $s$ be arbitrary scalar functions on $\Omega\times \mathbb R\times \mathbb R^n$ we obtain, by adding


% PDF page 376
\source{gilbarg-trudinger:page-0376}
\pagestyle{empty}


$\bar{v}(\omega(r+1)+s)$ times equation (15.17) to equation (15.21), the equation

\begin{equation}
\begin{aligned}
&-a^{ij}D_{ik}\bar{u}D_{jk}\bar{u}+\tfrac12a^{ij}D_{ij}\bar{v}
+\tfrac12\bigl(\psi' D_{p_i}a^{jk}D_{jk}\bar{u}+D_{p_i}b+\omega D_{p_i}\mathcal{E}\bigr)D_i\bar{v}\\
&\qquad+\psi'\bar{v}\bigl(\omega(\delta+r+1)a^{ij}+(\delta+s)a^{ij}\bigr)D_{ij}\bar{u}\\
&\qquad+\left\{\frac{\omega'}{\psi'}\mathcal{E}+\omega^2(\delta+r)\mathcal{E}
+\omega\bigl((\delta+s)\mathcal{E}+(\delta+r)b\bigr)+(\delta+s)b\right\}\bar{v}=0.
\end{aligned}
\tag{15.22}
\end{equation}

At this point it is convenient to write the principal coefficients $a^{ij}$ in the form

\begin{equation}
a^{ij}(x,z,p)=a^{ij}_{*}(x,z,p)+\tfrac12\bigl[p_ic_{j}(x,z,p)+c_i(x,z,p)p_j\bigr]
\tag{15.23}
\end{equation}

where $a^{ij}_{*},\,c_i\in C^1(\Omega\times\mathbb{R}\times(\mathbb{R}^n-\{0\}))$ and
$[a^{ij}_{*}]$ is a positive, symmetric matrix. Clearly the decomposition (15.2) with positive
$[a^{ij}]$ is a special case of (15.23). Furthermore the principal coefficients of any elliptic
operator $Q$ can be written in the form (15.23) by simply choosing $a^{ij}_{*}=a^{ij}$ and
$c_i=0$. Indeed for our intended application to uniformly elliptic equations there is nothing
to be gained by taking non-trivial $c_i$. However for the minimal surface operator, a decomposition
of the form (15.23) with the matrix $[a^{ij}_{*}]$ proportional to the identity matrix is crucial
for our derivation of global gradient bounds.

Returning to equation (15.22) we now substitute the relation (15.23) to obtain

\begin{equation}
\begin{aligned}
&-a^{ij}_{*}D_{ik}\bar{u}D_{jk}\bar{u}+\tfrac12a^{ij}D_{ij}\bar{v}
+\tfrac12\bigl(\psi' D_{p_i}a^{jk}D_{jk}\bar{u}-\psi'c_jD_{ij}\bar{u}\\
&\qquad+\bar{v}\bigl[\omega(\delta+r+1)+\delta+s+1\bigr]c_i+D_{p_i}b+\omega D_{p_i}\mathcal{E}\bigr)D_i\bar{v}\\
&\qquad+\psi'\bar{v}\bigl[\omega(\delta+r+1)a^{ij}+(\delta+s)a^{ij}_{*}\bigr]D_{ij}\bar{u}\\
&\qquad+\left\{\frac{\omega'}{\psi'}\mathcal{E}+\omega^2(\delta+r)\mathcal{E}
+\omega\bigl[(\delta+s)\mathcal{E}+(\delta+r)b\bigr]+(\delta+s)b\right\}\bar{v}=0.
\end{aligned}
\tag{15.24}
\end{equation}

By Cauchy's inequality (7.6) we can estimate

\[
\psi'\bar{v}\bigl[\omega(\delta+r+1)a^{ij}+(\delta+s)a^{ij}_{*}\bigr]D_{ij}\bar{u}
\]
\[
\le \lambda_{*}\sum\left|D_{ij}\bar{u}\right|^2+\frac{\bar{v}}{4\lambda_{*}}
\sum\left[\omega(\delta+r+1)+\delta+s\right]\left[a^{ij}_{*}\right]^2
\]
\[
\le a^{ij}_{*}D_{ik}\bar{u}D_{jk}\bar{u}+\frac{\bar{v}}{4\lambda_{*}}
\sum\left[\omega(\delta+r+1)+\delta+s\right]\left[a^{ij}_{*}\right]^2
\]

where $\lambda_{*}$ denotes the minimum eigenvalue of the matrix $[a^{ij}_{*}]$. Hence by substitution
into (15.24) we finally obtain the inequality

\begin{equation}
a^{ij}D_{ij}\bar{v}+B_iD_i\bar{v}+2G\mathcal{E}\bar{v}\ge 0
\tag{15.25}
\end{equation}

where the coefficients $B_i$ and $G$ are given by

\begin{equation}
B_i=\psi'\bigl(D_{p_i}a^{jk}D_{jk}\bar{u}-c_jD_{ij}\bar{u}\bigr)+\bar{v}\bigl[\omega(\delta+r+1)+\delta+s+1\bigr]c_i
+D_{p_i}b+\omega D_{p_i}\mathcal{E},
\tag{15.26}
\end{equation}

\[
G=\frac{\omega'}{\psi'}+\alpha\omega^2+\beta\omega+\gamma.
\]


% PDF page 377
\source{gilbarg-trudinger:page-0377}
\[
\begin{aligned}
\alpha &= \frac{1}{\mathcal{E}_*}\left((\delta+r)\mathcal{E} + \frac{|p|^2}{4\lambda_*}\sum\bigl((\delta+r+1)a_*^{ij}\bigr)^2\right),\\
\beta &= \frac{1}{\mathcal{E}_*}\left((\delta+s)\mathcal{E}+(\delta+r)b+\frac{|p|^2}{2\lambda_*}\bigl[(\delta+r+1)a_*^{ij}\bigr]\bigl[(\delta+s)a_*^{ij}\bigr]\right),\\
\gamma &= \frac{1}{\mathcal{E}_*}\left[\frac{|p|^2}{4\lambda_*}\sum\bigl((\delta+s)a_*^{ij}\bigr)^2+(\delta+s)b\right].
\end{aligned}
\]

In order to get a global gradient bound for the solution \(u\), we need to choose the functions \(\psi\), \(r\) and \(s\) in such a way that

\begin{equation}\tag{15.28}
G\le 0 \qquad \text{for } x\in \Omega,\qquad \text{and}\qquad |Du(x)|\ge L
\end{equation}

for sufficiently large \(L\). For, if (15.28) is satisfied, it follows from the weak maximum principle, Theorem 3.1, that

\[
\sup_{\Omega_L}\bar v=\sup_{\partial\Omega_L}\bar v,\qquad
(\Omega_L=\{x\in \Omega\mid |Du(x)|\ge L\}),
\]

and hence

\begin{equation}\tag{15.29}
\sup_{\Omega}|Du|\le \max\left\{\left[\frac{\max \psi'}{\min \psi'}\right]\sup_{\partial\Omega}|Du|,\,L\right\},
\end{equation}

We digress here momentarily to remark that the estimate (15.29) would remain valid if the solution \(u\) is only assumed to lie in \(C^2(\Omega)\). For, in this case the equation (15.18) would continue to hold in the weak form

\begin{equation}\tag{15.30}
\begin{aligned}
\int_{\Omega}\eta a^{ij}D_{ik}\bar u\,D_{jk}\bar u\,dx
&+\frac{1}{2}\int_{\Omega} a^{ij}D_i\bar v\,D_j\eta\,dx\\
+\frac{1}{2}\int_{\Omega}\left\{D_j a^{ij}D_i\bar v-\frac{1}{(\psi')^2}\bigl(\psi' D_p a^{jk}D_{jk}\bar u + D_p b + \omega D_p \mathcal{E}^{\circ}\bigr)D_i\bar v\right\}\eta\,dx\\
-\int_{\Omega}\left\{\psi'\,\delta a^{ij}D_{ij}\bar u+\delta b-\omega(b+\omega\mathcal{E}^{\circ})+\frac{\omega'}{\psi'}\,\mathcal{E}^{\circ}\right\}\bar v\eta\,dx=0
\end{aligned}
\end{equation}

for all \(\eta\ge 0\), \(\eta\in C_0^1(\Omega)\)

(see Sections 13.3 or 14.1). Applying the same argument as above for the case \(u\in C^3(\Omega)\) we deduce from (15.30) the weak form of inequality (15.25), namely

\begin{equation}\tag{15.31}
\int_{\Omega}\left\{a^{ij}D_i\bar v\,D_j\eta+(D_j a^{ij}-B_i)D_i\bar v\,\eta-2G\mathcal{E}^{\circ}\bar v\eta\right\}\,dx\le 0
\end{equation}

% PDF page 378
\source{gilbarg-trudinger:page-0378}
\pagestyle{empty}


for all non-negative $\eta \in C^{1}_{0}(\Omega)$. The estimate $(15.29)$ is then a consequence of the weak maximum principle, Theorem 8.1.

Let us now consider conditions on the coefficients of $Q$ which will ensure that inequality $(15.28)$ is satisfied for some $\psi$, $r$ and $s$. In order to guarantee that $\alpha$, $\beta$ and $\gamma$ are bounded from above, we impose the structure conditions

\begin{equation}
(15.32)\qquad (\delta+r+1)a^{ij}_{*},\,(\delta+s)a^{ij}_{*}=O\!\left[\frac{\sqrt{\lambda_{*}\mathcal{E}}}{|p|}\right]
\quad\text{as }|p|\to\infty
\end{equation}

\[
\delta\mathcal{E},\,\delta\mathcal{E},\,(\delta+r)b,\,(\delta+s)b \le O(\mathcal{E}) \quad \text{as }|p|\to\infty.
\]

Here the limit behavior is understood to be uniform for $(x,z)\in \Omega\times [m,M]$. We now define the numbers

\begin{equation}
(15.33)\qquad a,\,b,\,c = \overline{\lim}_{|p|\to\infty}\sup_{\Omega\times [m,M]}\ \alpha,\,\beta,\,\gamma
\end{equation}

so that inequality $(15.28)$ is implied by the Riccati differential inequality

\begin{equation}
(15.34)\qquad \chi'(z)+a\chi^{2}(z)+b\chi(z)+c+\varepsilon\le 0
\end{equation}

holding on the interval $[m,M]$ for some positive number $\varepsilon$. To get from the solution $\chi$ to our auxiliary function $\psi$ we use the relation $\chi=\omega\circ\psi^{-1}$. It turns out that inequality $(15.34)$ cannot be solved for arbitrary $a$, $b$ and $c$ unless $\operatorname{osc}_{\Omega}u=M-m$

is sufficiently small (see Problem 15.1). The inequality can be solved however when either $a\le 0$, $c\le 0$ or $b\le -2\sqrt{|ac|}$. Let us consider now the two important cases $a\le 0$, $c\le 0$; (the case $b\le -2\sqrt{|ac|}$ is left to the reader in Problem 15.2). We simplify our calculations somewhat by setting $\varphi(z)=\psi\circ\psi^{-1}(z)$. Then if $\chi$ satisfies inequality $(15.34)$ we can determine $\varphi$ through the relation

\[
\frac{\varphi'}{\varphi}=\chi.
\]

\begin{enumerate}
\item[(i)] $a\le 0$. If the strict inequality $a<0$ holds, the quadratic equation
\[
a\chi^{2}+b\chi+c+\varepsilon=0
\]
has a positive real root $\chi=\kappa$, if $\varepsilon$ is chosen so that $c+\varepsilon>0$, in which case inequality $(15.34)$ is solved by taking $\chi=\kappa$. Consequently we obtain $(\log \varphi)'=\kappa$ and hence we can choose
\[
\varphi(z)=e^{\kappa z}.
\]

If, on the other hand, $a=0$, the inequality $(15.34)$ can be solved by taking
\[
\chi(z)=\frac{|c+\varepsilon|}{|b|+\varepsilon}e^{2(|b|+\varepsilon)(M-z)}
\]
\end{enumerate}


% PDF page 379
\source{gilbarg-trudinger:page-0379}

in which case we can choose
\[
\varphi(z)=\exp\left\{-\frac{|c+\epsilon|}{2(|b|+\epsilon)^2}e^{2(|b|+\epsilon)(M-z)}\right\}.
\]

(ii) $c\leq 0$. If $c<0$ we can take $\varphi(z)=e^{\kappa z}$ for some $\kappa>0$ as in (i). If $c=0$, inequality (15.34) can be solved, for sufficiently small $\epsilon$, by taking
\[
\chi(c)=e^{A(m-z-1)} \qquad \text{where } A=|a|+|b|+1.
\]

In this case we can choose
\[
\varphi(z)=\exp\left\{-\frac{1}{A}e^{A(m-z-1)}\right\}.
\]

Note that in the cases considered above $\varphi$ is monotone increasing and hence in the estimate (15.29) we have $\max_{\Omega}\psi'=\varphi(M)$, $\min_{\Omega}\psi'=\varphi(m)$.
The development of this section has accordingly led us to the establishment of
the following theorem.

\begin{theorem}[Theorem 15.2]\label{gt:thm-15-2}\dcref{15.2}\group{ch15-global-interior-gradient}\source{gilbarg-trudinger:page-0379}
\textbf{Theorem 15.2.} \textit{Let $u\in C^2(\Omega)\cap C^1(\bar{\Omega})$ satisfy equation (15.1) in the bounded domain $\Omega$. Suppose that the operator $Q$ is elliptic in $\Omega$ and that there exist scalar multipliers $r$ and $s$ such that the structure conditions (15.32) are fulfilled together with either of the conditions $a\leq 0$ or $c\leq 0$ (the numbers $a$ and $c$ being defined by (15.27) and (15.33)). Then we have the estimate}

\begin{equation}
\tag{15.35}
\sup_{\Omega} |Du| \leq C
\end{equation}

\textit{where $C$ depends on the quantities in (15.32), $\operatorname{osc}_{\Omega} u$ and $\sup_{\partial\Omega} |Du|$.}
\end{theorem}

We illustrate the application of Theorem 15.2 by considering some important special cases.

(i) \textit{Uniformly elliptic equations.} If $Q$ is uniformly elliptic in $\Omega$, that is $a^{ij}=O(\lambda)$ as $|p|\to\infty$, and if also the function $b=O(\epsilon)=O(\lambda|p|^2)$ as $|p|\to\infty$, then conditions (15.32) with $a^{ij}=a^{ij}$, $c_i=0$, $r=s=0$ are often referred to as \textit{natural conditions} (see [LU 4]). If we restrict these conditions slightly by requiring that

\begin{equation}
\tag{15.36}
\delta a^{ij}=o(\lambda), \qquad \delta b=o(\lambda|p|^2) \qquad \text{as } |p|\to\infty
\end{equation}

then $c\leq 0$, and hence $C^2(\Omega)$ solutions of the equation $Qu=0$ satisfy an a priori global gradient bound. We shall show in the next section that the restriction (15.36) is unnecessary.

(ii) \textit{Prescribed mean curvature equation.} By writing equation (10.7) in the form

\begin{equation}
\tag{15.37}
\Delta u-\frac{D_i u\, D_j u}{(1+|Du|^2)}\,D_{ij}u-nH(x)\sqrt{1+|Du|^2}=0
\end{equation}

% PDF page 380
\source{gilbarg-trudinger:page-0380}

where $H \in C^{1}(\overline{\Omega})$ we can choose

\[
a_{*}^{ij}=\delta^{ij}, \qquad c_i=-\frac{p_i}{1+|p|^{2}}, \qquad r=-1 \quad \text{and} \quad s=0
\]

so that by calculation we have

\[
\alpha=-1+\frac{2}{1+|p|^{2}}, \qquad \beta=\frac{n H(x)\sqrt{1+|p|^{2}}}{|p|^{2}},
\]

\[
\gamma=-\frac{n p_i D_i H \, (1+|p|^{2})^{3/2}}{|p|^{4}}
\]

and hence

\[
a=-1, \qquad b=0, \qquad c=n \sup_{\Omega} |DH|.
\]

Consequently, a global gradient bound holds for $C^{2}(\Omega)$ solutions by virtue of the case $a<0$ of Theorem 15.2. In particular, we can deduce using (15.29) the estimate

\[
(15.38)\qquad \sup_{\Omega} |Du| \le c_1 + c_2 n \sup_{\Omega} |H| + c_3 \sup_{\partial\Omega} |Du| \exp\!\bigl(c_4 n \sup_{\Omega} |DH| \, \osc_{\Omega} u\bigr)
\]

where $c_1, c_2, c_3$ and $c_4$ are constants.

\[
\text{(iii)}\quad \textit{Equations where } a_{*}^{ij}=\delta^{ij}.
\text{ As shown in the preceding section, if equation}
\]
\[
(15.1) \text{ is the Euler-Lagrange equation of a multiple integral of the form (15.3) or}
\]
\[
\text{if the dimension } n=2,\text{ then we can take } a_{*}^{ij}=\delta^{ij}. \text{ Using (15.27) we obtain in these}
\]
\[
\text{cases}
\]

\[
\alpha=\frac{1}{\mathscr{E}}(\delta+r)\mathscr{E}+\frac{(r+1)^{2}|p|^{2}}{4}\frac{1}{\mathscr{E}},
\]

\[
(15.39)\qquad \beta=\frac{1}{\mathscr{E}}\bigl[(\delta+s)\mathscr{E}+(\delta+r)b\bigr]+\frac{(r+1)s|p|^{2}}{2}\frac{1}{\mathscr{E}},
\]

\[
\gamma=\frac{1}{\mathscr{E}}(\delta+s)b+\frac{s^{2}}{4}\frac{|p|^{2}}{\mathscr{E}}.
\]

Choosing the particular values $r=-1$, $s=0$ we have

\[
\alpha=\frac{1}{\mathscr{E}}(\delta-1)\mathscr{E},
\]

\[
\beta=\frac{1}{\mathscr{E}}\bigl[\delta\mathscr{E}+(\delta-1)b\bigr],
\]

\[
\gamma=\frac{1}{\mathscr{E}}\delta b.
\]

% PDF page 381
\source{gilbarg-trudinger:page-0381}

Note that the same formulae arise whenever

\[
a_{*}^{ij}(p)=a_{*}^{ij}(p/|p|)
\]

since in these cases we have $\delta a_{*}^{ij}=0$.

\bigskip

\section*{15.3. Interior Gradient Bounds}

An interior gradient bound for $C^{2}(\Omega)$ solutions of equation (15.1) can also be deduced from equation (15.24). Let $B=B_{R}(y)$ be a ball strictly lying in $\Omega$ and let $\eta$ be a function in $C^{2}(\overline{B})$ such that $0 \le \eta \le 1$ in $\overline{B}$, $\eta=0$ on $\partial B$, $\eta(y)=1$ and $\eta>0$ in $B$. A typical example of such a function, which we shall use below, is given by

\[
(15.40)\qquad \eta(x)=\left(1-\frac{|x-y|^{2}}{R^{2}}\right)^{\beta},\qquad \beta \ge 1.
\]

We now consider the function $w$ defined by

\[
w=\eta \bar{v}
\]

Clearly $w \in C^{1}(\overline{B})$, $w(y)=\bar{v}(y)$, $w=0$ on $\partial B$, and the derivatives of $w$ are given by

\[
D_{i}w=\eta D_{i}\bar{v}+\bar{v}D_{i}\eta
\]
\[
D_{ij}w=\eta D_{ij}\bar{v}+D_{i}\bar{v}D_{j}\eta+D_{i}\eta D_{j}\bar{v}+\bar{v}D_{ij}\eta.
\]

Note that to guarantee the existence of the second derivatives $D_{ij}w$ we require that $u \in C^{3}(\Omega)$; this restriction can however be avoided by utilizing the weak form of equation (15.24). Multiplying equation (15.24) by $\eta$ and substituting for $\bar{v}$, we now obtain the equation

\[
-\eta a_{*}^{ij}D_{ik}\bar{u}D_{jk}w+\frac{1}{2}a^{ij}D_{ij}w+\frac{1}{2}\left(B_{i}-\frac{2}{\eta}a^{ij}D_{j}\eta\right)D_{i}w
\]
\[
+\psi' w[\omega(\delta+r+1)+(\delta+s)]a_{*}^{ij}D_{ij}\bar{u}
\]
\[
+\left\{\frac{\omega'}{\psi'}\mathscr{E}+\omega^{2}(\delta+r)\mathscr{E}+\omega[(\delta+s)\mathscr{E}+(\delta+r)b]+(\delta+s)b\right.
\]
\[
\qquad\qquad\left.+\frac{1}{\eta^{2}}a^{ij}D_{i}\eta D_{j}\eta-\frac{1}{2\eta}a^{ij}D_{ij}\eta-\frac{1}{2\eta}B_{i}D_{i}\eta\right\}w=0,
\]

where $B_{i}$ is given by (15.26). At this point we can introduce further scalar multipliers $t_{i}$, $i=1,\ldots,n$. For, if $t_{i}$, $i=1,\ldots,n$, are arbitrary scalar functions on $\Omega \times \mathbb{R} \times \mathbb{R}^{n}$ we can write, using equation (15.17) and setting $c_{i}=D_{p_{i}},$

% PDF page 382
\source{gilbarg-trudinger:page-0382}

\[
B_i=\psi'\bigl[(c_i+t_i)a^{jk}D_{jk}\bar u-c_jD_{ij}\bar u\bigr]
   +v[\omega(\delta+r+1)+(\delta+s+1)]c_i
   +(c_i+t_i)b+\omega(c_i+t_i)\mathcal E
\]
\[
=\psi'(c_i+t_i)a^{jk}D_{jk}\bar u-\frac{v}{2\eta}D_i\eta(c_i+t_i)c_j
  +\frac{(\psi')^2}{2\eta}(c_i+t_i)c_jD_jw
\]
\[
\qquad +v[\omega(\delta+r+1)+(\delta+s+1)]c_i+(c_i+t_i)b+\omega(c_i+t_i)\mathcal E.
\]

Hence by substitution into equation (15.41), we have

\begin{equation}
\begin{aligned}
&-\eta a_*^{ij}D_i\bar u\,D_j\bar u+\frac12 a^{ij}D_{ij}w+\frac12 \widetilde B_iD_iw\\
&\quad +\psi'w[\omega(\delta+r+1)+(\delta+s)+(\tilde\delta+t)]a_*^{ij}D_{ij}\bar u\\
&\quad +\Bigl\{\frac{\omega'}{\psi'}\mathcal E+\omega^2(\delta+r)\mathcal E+\omega[(\delta+\tilde\delta+s+t)\mathcal E+(\delta+r)b\\
&\qquad\qquad -\frac{v}{2\eta}D_i\eta(\delta+r+1)c_i]+(\delta+\tilde\delta+s+t)b\\
&\qquad\qquad -\frac{v}{2\eta}D_i\eta(\delta+\tilde\delta+s+t+1)c_i\\
&\qquad\qquad +\frac{1}{\eta^2}a^{ij}D_i\eta D_j\eta-\frac{1}{2\eta}a^{ij}D_{ij}\eta\Bigr\}w=0
\end{aligned}
\tag{15.42}
\end{equation}

where

\begin{equation}
\tilde\delta=-\frac{1}{2\eta}D_i\eta\,c_i,
\qquad
t=-\frac{t_i}{2\eta}D_i\eta
\tag{15.43}
\end{equation}

and

\[
\widetilde B_i=B_i-\frac{2}{\eta}a^{ij}D_j\eta-\frac{v}{2\eta}D_j\eta(\delta_j+t_j)c_i.
\]

We then obtain, instead of inequality (15.25), the inequality

\begin{equation}
a^{ij}D_{ij}w+\widetilde B_iD_iw+2\widetilde G\mathcal Ew\ge 0
\tag{15.44}
\end{equation}

where $\widetilde G$ is given by

\begin{equation}
\widetilde G=\frac{\omega'}{\psi'}+\tilde\alpha\omega^2+\beta\omega+\tilde\gamma,
\tag{15.45}
\end{equation}

% PDF page 383
\source{gilbarg-trudinger:page-0383}

and

\[
\tilde{a}=\alpha,
\]

\[
\tilde{\beta}=\beta+\frac{1}{\mathscr{E}}\left\{-\frac{|p|^{2}}{2\eta}D_{i}\eta(\delta+r+1)c_{i}+(\delta+t)\mathscr{E}
\qquad\qquad\qquad\qquad+\frac{|p|^{2}}{2\lambda_{*}}\left[(\delta+r+1)a_{*}^{ij}\right]\left[(\tilde{\delta}+t)a_{*}^{ij}\right]\right\},
\]

\begin{equation}
\tag{15.46}
\tilde{\gamma}=\gamma+\frac{1}{\mathscr{E}}\left\{\frac{|p|^{2}}{4\lambda_{*}}\sum \left[(\delta+t)a_{*}^{ij}\right]^{2}+(\delta+t)b\right.
\qquad\qquad-\frac{v}{2\eta}D_{i}\eta(\delta+\tilde{\delta}+s+t+1)c_{i}+\frac{1}{\eta^{2}}a^{ij}D_{i}\eta D_{j}\eta
\qquad\qquad-\frac{1}{2\eta}a^{ij}D_{ij}\eta+\frac{|p|^{2}}{2\lambda_{*}}\left[(\delta+s)a_{*}^{ij}\right]\left[(\tilde{\delta}+t)a_{*}^{ij}\right]\Biggr\}.
\end{equation}

Unless there exist special relationships between the function $\eta$ and the coefficients
of $Q$, we need to add more structure conditions to ensure that the coefficients $\tilde{\beta}$
and $\tilde{\gamma}$ behave similarly to $\beta$ and $\gamma$. Let us therefore assume, in addition to the
structure conditions (15.32), the conditions

\[
|p|^{\theta}(\partial_{k}+t_{k})a_{*}^{ij}=O\!\left(\sqrt{\lambda_{*}\mathscr{E}}/|p|\right)
\qquad \text{as } |p|\to\infty,
\]

\begin{equation}
\tag{15.47}
|p|^{2\theta}\lambda_{*},\,|p|^{\theta}(\partial_{i}+t_{i})\mathscr{E},\,|p|^{\theta}(\partial_{i}+t_{i})b=O(\mathscr{E})
\qquad \text{as } |p|\to\infty,
\end{equation}

\[
|p|^{\theta}(\tilde{\delta}+r+1)c_{i},\,|p|^{\theta}(\delta+s+1)c_{i},\,|p|^{2\theta}(\partial_{i}+t_{i})c_{i}=O(\mathscr{E}/|p|^{2})
\]

\[
\qquad\qquad\text{as } |p|\to\infty,
\]

with $i,j,k=1,\ldots,n$, for some $\theta>0$. As in (15.32), the limit behavior is understood
to be uniform for $(x,z)\in(m,M)$.

Using (15.47) and the function $\eta$ given by (15.40), with $\beta=2/\theta$, we can then
estimate (for sufficiently large $|Du|$)

\[
|\beta-\tilde{\beta}|\leq C\frac{|D\eta|}{\eta|Du|^{\theta}}\leq \frac{C\beta}{Rw^{1/\beta}},
\]

\[
|\gamma-\tilde{\gamma}|\leq C\left\{\frac{|D\eta|^{2}}{\eta^{2}|Du|^{2\theta}}+\frac{|D\eta|}{\eta|Du|^{\theta}}\right\}\leq
\left\{\frac{\beta^{2}}{R^{2}w^{2/\beta}}+\frac{\beta}{Rw^{1/\beta}}\right\},
\]

where $C$ depends on $n$ and the quantities in (15.32) and (15.47). Consequently,
applying the considerations of the preceding section we obtain, instead of (15.29),
the estimate

\begin{equation}
\tag{15.48}
|Du(y)|\leq C\left(1+R^{-1/\theta}\right),
\end{equation}

% PDF page 384
\source{gilbarg-trudinger:page-0384}
\pagestyle{empty}

provided either $a \leq 0$, $c \leq 0$, $b \leq -2\sqrt{ac}$, or $\osc_{B_R(y)}u$ is less than some constant depending on $a,b,c$. Moreover, we can clearly replace the ball $B = B_R(y)$ by any ball intersecting $\Omega$, and therefore obtain, for any $y \in \Omega$, the same estimate with $C$ depending in addition on $\sup_{B\cap \Omega}|Du|$. Note also that by invoking (15.30), rather than (15.18), the above considerations are also applicable to $C^2(\Omega)$ solutions. Therefore, we have the following estimate.

\textbf{Theorem 15.3.} \textit{Let $u \in C^2(\Omega)$ satisfy (15.1) in the domain $\Omega$. Suppose that the operator $Q$ is elliptic in $\Omega$ and that there exist scalar multipliers $r,s,t_i$, $i=1,\ldots,n$ such that the structure conditions (15.32) and (15.47) are fulfilled. Then, if any of the conditions $a \leq 0$, $c \leq 0$, $b \leq -2\sqrt{ac}$ hold, (the numbers $a,b,c$ being defined by (15.27) and (15.33)), we have the estimate}

\begin{equation}
(15.49)\qquad |Du(y)| \leq C(1 + R^{-1/\theta})
\end{equation}

\textit{for any $y \in \Omega$ and $B = B_R(y)$, where $C$ depends on the quantities in (15.32), (15.47), $\osc_{B\cap \Omega}u$ and $\sup_{B\cap \Omega}|Du|$. If $a,b$ and $c$ are arbitrary, the estimate (15.49) continues to hold for sufficiently small $R$ depending on $a,b,c$ and the modulus of continuity of $u$ at $y$.}

For uniformly elliptic equations, under the natural conditions, estimates for the moduli of continuity of solutions (and hence gradient estimates) follow from the considerations of Section 9.7. In fact, we have the following estimate which, in some sense, extends Corollary 9.25.

\textbf{Lemma 15.4.} \textit{Let $u \in W^{2,n}(\Omega)$ satisfy the inequality}

\begin{equation}
(15.50)\qquad |Lu| \leq \lambda(\mu_0|Du|^2 + f)
\end{equation}

\textit{in $\Omega$, where the operator $L$, given by $Lu = a^{ij}D_{ij}u$, satisfies (9.47), $\mu_0 \in \mathbb{R}$ and $f \in L^n(\Omega)$. Then, for any ball $B_0 = B_{R_0}(y) \subset \Omega$ and $R \leq R_0$, we have}

\begin{equation}
(15.51)\qquad \osc_{B_R(y)} u \leq C\left(\frac{R}{R_0}\right)^\alpha \left(\osc_{B_0}u + R\|f\|_{n;\Omega}\right)
\end{equation}

\textit{where $C$ and $\alpha$ depend only on $n$, $\Lambda/\lambda$ and $\mu_0 M$; $M = \|u\|_{0;\Omega}$.}

\textit{Proof.} \ Let us suppose first that $u \geq 0$ satisfies

\begin{equation}
(15.52)\qquad Lu \leq \lambda(\mu_0|Du|^2 + f)
\end{equation}

\textit{in $\Omega$, where $\mu_0 > 0$, $f \geq 0$, and set}

\[
w = \frac{1}{\mu_0}\left(1 - e^{-\mu_0 u}\right)
\]


% PDF page 385
\source{gilbarg-trudinger:page-0385}

It follows then that

\[
Lw \leq f
\]

in $\Omega$, so that the weak Harnack inequality (Theorem 9.22) is applicable to $w$. But since

\[
w \leq u \leq e^{\mu_0 M}w,
\]

the weak Harnack inequality (9.48) is satisfied also by $u$, with constant $C$ depending in addition on $\mu_0 M$. The H\"older estimate (15.51) then follows.

Combining Theorem 15.3 and Lemma 15.4, we conclude both interior and global gradient bounds under the conditions

\begin{equation}
\Lambda,(\delta + r + 1)a^{ij},(\delta + s)a^{ij},|p|^\theta(\partial_k + t_k)a^{ij} = O(\lambda),
\end{equation}
\[
b,|p|^\theta(\partial_i + t_i)b = O(\lambda|p|^2),(\delta + r)b,(\delta + s)b \leq O(\lambda|p|^2),
\]

as $|p| \to \infty$, for some $\theta > 0$.

\textbf{Theorem 15.5.} \emph{Let $u \in C^2(\Omega)$ satisfy (15.1) in the domain $\Omega$ and suppose that the operator $Q$ is elliptic in $\Omega$ and satisfies the structure conditions (15.53), with multipliers $r, s, t_i, i = 1,\ldots,n$. Then, for any point $y \in \Omega$, we have}

\begin{equation}
|Du(y)| \leq C\left(1 + d_y^{-1/\theta}\right)
\end{equation}

\emph{where $C$ depends on $n$, the quantities in (15.53), $\|u\|_{0;\Omega}$, and $d_y = \operatorname{dist}(y,\partial \Omega)$. If $u \in C^2(\Omega)\cap C^1(\bar{\Omega})$, we also have}

\begin{equation}
|Du(y)| \leq C
\end{equation}

\emph{where $C$ depends in addition on $\sup_{\partial\Omega}|Du|$.}

\bigskip

\section*{15.4. Equations in Divergence Form}

Let $u \in C^2(\Omega)$ satisfy the divergence form equation

\begin{equation}
Qu = \divop A(x,u,Du) + B(x,u,Du) = 0
\end{equation}

in the domain $\Omega$. Then it was shown in Section 13.1 that the derivatives $D_k u$, $k = 1,\ldots,n$ satisfy the linear divergence form equations

\begin{equation}
\int_\Omega \left(\bar a^{ij}D_jD_k u + f_k^i\right)D_i\zeta\,dx = 0 \qquad \forall \zeta \in C_0^1(\Omega),
\end{equation}


% PDF page 386
\source{gilbarg-trudinger:page-0386}


where

\[
\tag{15.58}
\bar{a}^{ij}(x)=D_{p_j}A^i(x,u(x),Du(x)),
\qquad
f^i_k(x)=\delta_k A^i(x,u(x),Du(x))+\delta^{ik}B(x,u(x),Du(x)),\qquad \delta_k=p_k D_z+D_{x_k}.
\]

In order to proceed further we shall assume that $Q$ is elliptic in the sense that

\[
\bar{a}^{ij}(x,z,p)\xi_i\xi_j=D_{p_j}A^i(x,z,p)\xi_i\xi_j\ge \nu(|z|)(1+|p|)^\tau |\xi|^2
\tag{15.59}
\]

for all $\xi\in \mathbb{R}^n$, $(x,z,p)\in \Omega\times \mathbb{R}\times \mathbb{R}^n$, where $\tau$ is some real number and $\nu$ is a positive, non-increasing function on $\mathbb{R}$. A global gradient bound for non-uniformly elliptic $Q$ can then be derived under the additional structure conditions

\[
|p|\,D_z A,\ D_x A,\ B=o(|p|)(1+|p|)^\tau \qquad \text{as } |p|\to\infty.
\tag{15.60}
\]

To show this we put $M=\sup_{\Omega}|u|$, $M_1=\sup_{\Omega}|Du|$ and apply the maximum principle, Theorem 8.16, to equation (15.57) in the domain $\tilde{\Omega}=\{x\in \Omega\mid |Du(x)|>\frac{M_1}{2\sqrt{n}}\}$. We obtain thus

\[
\sup_{\tilde{\Omega}} |D_k u|\le \sup_{\partial\tilde{\Omega}} |D_k u|+C(2\sqrt{n})^\tau \|(1+|Du|)^{-\tau}f^i_k\|_{q}
\tag{15.61}
\]

for $q>n$, $k=1,\dots,n$, where $C=C(n,\nu(M),q,|\Omega|)$. Taking $q=\infty$ and using conditions (15.60), we therefore have

\[
M_1=\sup_{\Omega}|Du|\le C\bigl(\sup_{\partial\Omega}|Du|+\sigma(M_1)\bigr)
\tag{15.62}
\]

where $\sigma(t)=o(t)$ as $t\to\infty$. Consequently, an apriori estimate for $M_1$ follows.

\medskip

\noindent\textbf{Theorem 15.6.} \textit{Let $u\in C^2(\Omega)$ satisfy equation (15.56) in the bounded domain $\Omega$ and suppose that the structure conditions (15.59), (15.60) are fulfilled. Then we have the estimate}

\[
\sup_{\Omega}|Du|\le C\bigl(1+\sup_{\partial\Omega}|Du|\bigr)
\tag{15.63}
\]

\textit{where $C$ depends on $n$, $\tau$, $\nu(M)$ and the quantities in (15.60).}

Condition (15.60) requires that the coefficient $b$ in equation (15.1) satisfies $b=o(\lambda|p|)$ as $|p|\to\infty$. In order to relax this condition we must either impose additional conditions on the coefficient matrix $D_pA$ such as uniform ellipticity or assume that the solution $u$ vanishes on $\partial\Omega$; (see Problems 15.4, 15.5). Note that the existence of a global gradient bound is reduced, through the estimate (15.61), to the existence of suitable integral estimates for the functions $(1+|Du|)^{-\tau}f^i_k$. Instead of pursuing global bounds any further, at this stage, we shall now turn to a

% PDF page 387
\source{gilbarg-trudinger:page-0387}

consideration of interior gradient estimates for uniformly elliptic equations. Let
us assume that the operator $Q$ is uniformly elliptic in $\Omega$ in the sense that

\begin{equation}
|D_p A(x,z,p)| \le \mu(|z|)(1+|p|)^\tau
\tag{15.64}
\end{equation}

for all $(x,z,p)\in \Omega\times \mathbb{R}\times \mathbb{R}^n$, where $\mu$ is a positive, non-decreasing function on $\mathbb{R}$.
Henceforth we shall also assume that $\tau>-1$, in which case conditions (15.59),
(15.64) imply respectively the inequalities

\[
p\cdot A(x,z,p)-p\cdot A(x,z,0)\ge \frac{\nu}{\tau+1}|p|^{\tau+2}.
\]

\begin{equation}
|A^i(x,z,p)-A^i(x,z,0)|\le \frac{\mu}{\tau+1}(1+|p|)^{\tau+1},\qquad i=1,\dots,n.
\tag{15.65}
\end{equation}

Finally we shall take, in place of (15.60), the more general condition

\begin{equation}
g(x,z,p)=(1+|p|)\,|D_z A|+|D_x A|+|B|\le \mu(|z|)(1+|p|)^{\tau+2}
\tag{15.66}
\end{equation}

for all $(x,z,p)\in \Omega\times \mathbb{R}\times \mathbb{R}^n$. Conditions (15.59), (15.64), (15.66) can be regarded
as natural for divergence structure operators. The derivation of an apriori interior
gradient bound under these conditions is accomplished in three stages:

\medskip

\noindent\textit{(i) Reduction to an $L^p$ estimate.} We replace the test function $\zeta$ in (15.57) by
$\zeta D_k u$ and sum the resulting equations over $k$. Setting $v=|Du|^2$ we obtain thus

\[
\int_{\Omega} \zeta \tilde{a}^{ij}D_i u D_j u\,dx+\int_{\Omega}\left(\frac12 \tilde{a}^{ij}D_j v + D_k u\,f_k^i\right)D_i\zeta\,dx+\int_{\Omega}\zeta\,f_k^i D_{ik}u\,dx=0.
\]

Hence by Young's inequality (7.6) we have

\begin{equation}
\int_{\Omega}\left(\tilde{a}^{ij}D_j v+2D_k u\,f_k^i\right)D_i\zeta\,dx\le \frac12\int_{\Omega}\zeta\lambda^{-1}\sum (f_k^i)^2\,dx
\tag{15.67}
\end{equation}

for all non-negative $\zeta\in C_0^1(\Omega)$. Therefore, setting

\[
\bar{v}=\int_0^v (1+\sqrt{t})^\tau\,dt,
\]

we may write (15.67) as

\begin{equation}
\int_{\Omega}\left(\tilde{a}^{ij}D_j \bar{v}+2D_k u\,f_k^i\right)D_i\zeta\le \frac12\int_{\Omega}\zeta\lambda^{-1}\sum (f_k^i)^2\,dx,
\tag{15.68}
\end{equation}

where

\[
\tilde{a}^{ij}=\bar{a}^{ij}(1+|Du|)^{-\tau}.
\]

% PDF page 388
\source{gilbarg-trudinger:page-0388}

The interior estimate for weak subsolutions of linear equations, Theorem 8.17, is thus applicable to inequality (15.68). Hence, using (15.66), we have for any ball $B_{2R}=B_{2R}(y)\subset \Omega$ and $q>n$, the estimate
\[
\sup_{B_R(y)} \bar v \le C\left\{R^{-n/2}\|\bar v\|_{L^2(B_{2R})}+\|(1+|Du|)^{\tau+4}\|_{L^q(B_{2R})}\right\}
\]
where $C=C(n,\nu(M),\mu(M),\tau,q,\operatorname{diam}\Omega)$, $M=\sup_{B_{2R}(y)}|u|$. Consequently for sufficiently large $p$ we have
\begin{equation}
\sup_{B_R(y)} v \le C(n,\nu(M),\mu(M),\tau,\operatorname{diam}\Omega,R^{-n}\int_{B_{2R}(y)} v^p\,dx).
\tag{15.69}
\end{equation}

\noindent (ii) \textit{Reduction to a Hölder estimate.} We now need to utilize the weak form of equation (15.57), viz.
\begin{equation}
\mathcal{Q}(u,\varphi)=\int_{\Omega}\bigl(A^i(x,u,Du)D_i\varphi-B(x,u,Du)\varphi\bigr)\,dx=0 \qquad \forall \varphi\in C_0^1(\Omega).
\tag{15.70}
\end{equation}

From (15.65) we see that the function $A$ satisfies inequalities
\begin{equation}
|A(x,z,p)|\le \mu_1(|z|)(1+|p|)^{\tau+1}
\tag{15.71}
\end{equation}
\[
p\cdot A(x,z,p)\ge \nu_1(|z|)|p|^{\tau+2}-\mu_1(|z|)
\]
for $(x,z,p)\in \Omega\times \mathbb{R}\times \mathbb{R}^n$, where $\mu_1$ and $\nu_1$ are respectively positive non-decreasing and positive non-increasing functions depending on $\tau,\mu,\nu$ and $\sup_{\Omega}|A(x,u,0)|$.

Let us substitute into (15.70) the test function
\[
\varphi=\eta^2[u-u(y)]
\]
where $\eta\in C_0^1(B_{2R})$, $B_{2R}=B_{2R}(y)\subset\Omega$. Using (15.66) and (15.71) we obtain thus
\[
\nu_1\int_{\Omega}\eta^2|Du|^{\tau+2}\,dx\le \mu_1\int_{\Omega}\eta^2\,dx
\]
\[
+\mu\int_{\Omega}\eta^2|u(x)-u(y)|(1+|Du|)^{\tau+2}\,dx
\]
\[
+2\mu_1\int_{\Omega}|\eta D\eta|\,|u(x)-u(y)|(1+|Du|)^{\tau+1}\,dx
\]
\[
\le \mu_1\int_{\Omega}\eta^2\,dx+4(\mu+\mu_1)2^\tau\omega(R)\int_{\Omega}\eta^2(1+|Du|^{\tau+2})\,dx
\]
\[
+\mu_1\omega(R)\int_{\Omega}(1+|Du|)^{\tau}|D\eta|^2\,dx
\]

% PDF page 389
\source{gilbarg-trudinger:page-0389}
where
\[
\omega(R)=\sup_{B_{2R}} |u(x)-u(y)|.
\]
Hence, if $R$ is chosen small enough to ensure that
\[
\omega(R)\leq \frac{\nu_1 2^{-\tau}}{8(\mu+\mu_1)},
\]
we have
\[
\int_{\Omega} \eta^2 |Du|^{\tau+2}\,dx \leq C \int_{\Omega} \bigl(\eta^2+\omega(R)(1+|Du|)^{\tau}|D\eta|^2\bigr)\,dx
\tag{15.72}
\]
where $C=C(\mu,\mu_1,\nu_1,\tau)$. Let us now replace the function $\eta$ in (15.72) by
\[
\eta v^{(\beta+1)/2}, \qquad \beta>0,
\]
to obtain the estimate
\[
\int_{\Omega} \eta^2 |Du|^{2\beta+\tau+4}\,dx \leq C \int_{\Omega} \Bigl\{ |Du|^{2(\beta+1)}\bigl(\eta^2+\omega(R)(1+|Du|)^{\tau}|D\eta|^2\bigr)
\qquad\qquad +(\beta+1)^2\omega(R)(1+|Du|)^{\tau}\eta^2 v^{\beta-1}|Dv|^2 \Bigr\}\,dx.
\tag{15.73}
\]
To estimate the last term in the above inequality we choose
\[
\zeta=\eta^2 v^{\beta}
\]
in inequality (15.67). We obtain thus, using conditions (15.59), (15.64) and (15.66),
\[
\beta v \int_{\Omega} (1+|Du|)^{\tau}\eta^2 v^{\beta-1}|Dv|^2\,dx
\]
\[
\leq C \int_{\Omega} \Bigl\{(1+|Du|)^{\tau}\eta v^{\beta}|D\eta|\,|Dv|+(1+|Du|)^{\tau+3}
\]
\[
\qquad\qquad \times \bigl(\eta v^{\beta}|D\eta|+\beta\eta^2 v^{\beta-1}|Dv|\bigr)+\eta^2(1+|Du|)^{\tau+4}v^{\beta}\Bigr\}\,dx
\]
where $C=C(\mu,\nu)$. Hence, by Young's inequality (7.6),
\[
\int_{\Omega} (1+|Du|)^{\tau}\eta^2 v^{\beta-1}|Dv|^2\,dx
\]
\[
\leq C \int_{\Omega} (1+|Du|)^{2\beta+\tau+2}\bigl(\eta^2(1+|Du|)^2+|D\eta|^2\bigr)\,dx
\]

% PDF page 390
\source{gilbarg-trudinger:page-0390}

where $C=C(\mu,\nu,\beta)$. Consequently if $\omega(R)$ is sufficiently small we obtain, by substitution into (15.73),
\[
\int_{\Omega}\eta^2(1+|Du|)^{2\beta+\tau+4}\,dx
\leq C\int_{\Omega}\left(\eta^2(1+|Du|)^{2(\beta+1)}+|D\eta|^2(1+|Du|)^{2\beta+\tau+2}\right)\,dx
\]
where $C=C(\mu,\nu,\beta,\tau)$. Replacing $\eta$ by $\eta^{(2\beta+\tau+4)/2}$ and using Young's inequality (7.6), we then obtain
\[
\int_{\Omega}[\eta(1+|Du|)]^{2\beta+\tau+4}\,dx\leq C
\]
where $C=C(\mu,\nu,\mu_1,\nu_1,\beta,\tau,\sup |D\eta|)$. In particular, if $\eta=1$ on $B_R(y)$ and $|D\eta|\leq 2/R$, it follows that, for any $p\geq 1$,
\begin{equation}\tag{15.74}
\int_{B_R(y)}(1+|Du|)^p\,dx\leq C
\end{equation}
where $C=C(\mu,\nu,\mu_1,\nu_1,p,\tau,R^{-1})$. Combining the estimates (15.69) and (15.74) we obtain for any ball $B_0=B_{R_0}(y)\subset\Omega$ and $0<\alpha<1$ the estimate
\begin{equation}\tag{15.75}
|Du(y)|\leq C
\end{equation}
where $C=C(n,\tau,\nu,\mu,\nu_1,\mu_1,\alpha,[u]_{\alpha,y})$ and
\[
[u]_{\alpha,y}=\sup_{B_0}\frac{|u(x)-u(y)|}{|x-y|^\alpha};
\]
that is, the derivation of an interior gradient bound is reduced to the existence of an interior H\"older estimate for $u$. We note here that an estimate of the modulus of continuity of $u$ would in fact suffice to complete the proof. Moreover, if the structure conditions (15.66) were strengthened so that $g=o(|p|^{t+2})$ as $|p|\to\infty$, the above considerations yield an interior gradient bound independent of the modulus of continuity of $u$. In this case we can also weaken the uniform ellipticity of $Q$ to allow $D_{p_j}A^i=O(|p|^{\tau+\sigma})$ as $|p|\to\infty$, $i,j=1,\ldots,n$, where $\sigma<1$ (see Problem 15.6).

\noindent (iii) \textit{A H\"older estimate for weak solutions of equation (15.66).} Let us write inequalities (15.71) together with the condition on $B$ in (15.66) in the form
\begin{equation}\tag{15.76}
\begin{aligned}
|A(x,z,p)|&\leq a_0|p|^{m-1}+\chi^{m-1} \\
p\cdot A(x,z,p)&\geq \nu_0|p|^m-\chi^m \\
|B(x,z,p)|&\leq b_0|p|^m+\chi^m
\end{aligned}
\end{equation}


% PDF page 391
\source{gilbarg-trudinger:page-0391}
\[
\text{where } m=\tau+2>1 \text{ and } a_0,\ b_0,\ v_0 \text{ and } \chi \text{ are positive constants depending possibly}
\]
\[
\text{on } M=\sup_{\Omega}|u|.\ \text{Then we have the following Hölder estimate.}
\]

\textbf{Theorem 15.7.} \emph{Let $u\in C^1(\Omega)$ be a weak solution of equation (15.56) in the domain
$\Omega$ and suppose that $\Omega$ satisfies the structure conditions (15.76). Then for any ball
$B_0=B_{R_0}(y)\subset\Omega$ and $R\leq R_0$ we have the estimate}

\begin{equation}
\tag{15.77}
\osc_{B_R(y)} u \leq C(1+R_0^{-\alpha}M_0)R^\alpha
\end{equation}

\emph{where $C=C(n,a_0,b_0,v_0,\chi,R_0,m,M_0)$ and $\alpha=\alpha(n,a_0,b_0,v_0,m,M_0)$ are positive
constants and $M_0=\sup_{B_0}|u|$.}

\textit{Proof.} Theorem 15.7 can be derived by essentially the same method as used for
Theorem 8.22 or alternatively by following the method described in Problem 8.6.
We shall not present all the details here. The essential ingredient in our proof of
Theorem 8.22 was the weak Harnack inequality Theorem 8.18. Setting
\[
k=\chi(R+R^{m/(m-1)})
\]
we obtain for non-negative weak supersolutions of equation (15.56) the analogous
weak Harnack inequality
\begin{equation}
\tag{15.78}
R^{-n/p}\|u\|_{L^p(B_{2R}(y))}\leq C\bigl(\inf_{B_R(y)} u+k\bigr)
\end{equation}
where $C=C(n,a_0,b_0,v_0,m,M_0,p)$ and $1\leq p<n/(n-m)^+$. The proof of (15.78)
can be modelled on that of Theorem 8.18 provided that in place of (8.48), we take
as test functions, in the inequality $\mathcal{Q}(u,\varphi)\geq 0$,
\[
\varphi=\eta^m \bar{u}^{\beta} \exp(-b_0\bar{u}), \qquad \beta<0,
\]
and use the Sobolev inequality (7.26) with $p=m$ instead of $p=2$. The passage
from the weak Harnack inequality to the Hölder estimate (15.77) can be accomplished
according to the proof of Theorem 8.22. Note that the case where $m>n$
can be handled directly from the Sobolev inequality, Theorem 7.17. \hfill$\square$

Combining Theorem 15.7 with our previous estimate (15.75) we finally arrive
at the following interior gradient estimate.

\textbf{Theorem 15.8.} \emph{Let $u\in C^2(\Omega)$ satisfy equation (15.57) in the domain $\Omega$ and suppose
that the structure conditions (15.60), (15.64) and (15.66) are fulfilled with $\tau>-1$.
Then we have the estimate}

\begin{equation}
\tag{15.79}
|Du(y)|\leq C
\end{equation}

% PDF page 392
\source{gilbarg-trudinger:page-0392}

for any $y \in \Omega$, where $C$ depends on $n$, $\tau$, $\nu(M_0)$, $\mu(M_0)$, $\sup_{\Omega} |A(x,u,0)|$, $M_0$, $M_0/d$ and $M_0 = \sup_{B_d(y)} |u|$, $d = \operatorname{dist}(y,\partial\Omega)$.

By applying an argument similar to the proof of Theorem 8.29, we can conclude from the interior estimate (15.79) (and the boundary Lipschitz estimate, Theorem 14.1) the following global estimate.

\textbf{Theorem 15.9.} \emph{Let $u \in C^2(\Omega) \cap C^0(\overline{\Omega})$ satisfy equation (15.57) in the bounded domain $\Omega$ and suppose that the structure conditions (15.60), (15.64) and (15.66) are fulfilled with $\tau > -1$. Assume also that $\Omega$ satisfies a uniform exterior sphere condition and that $u = \varphi$ on $\partial\Omega$ for $\varphi \in C^2(\overline{\Omega})$. Then we have the estimate}

\[
(15.80)\qquad \sup_{\Omega} |Du| \leq C
\]

\emph{where $C$ depends on $n$, $\tau$, $\nu(M)$, $\mu(M)$, $\sup_{\Omega} |A(x,u,0)|$, $\partial\Omega$, $|\varphi|_2$ and $M = \sup_{\Omega} |u|$.}

\medskip

\noindent\textbf{15.5. Selected Existence Theorems}

It is not feasible to present here a comprehensive account of existence theorems for the classical Dirichlet problem that follow by combination of the results of Chapters 10, 13, 14 and 15. Instead we shall present a selection of results which hopefully will serve as an illustration of the scope of the theory.

\noindent (i) \emph{Uniformly elliptic equations in the general form} (15.1). We assume the structure conditions

\[
a^{ij},\,\delta a^{ij} = O(\lambda),
\qquad
\delta a^{ij} = o(\lambda),
\]
\[
(15.81)\qquad b = O(\lambda |p|^2)
\]
\[
\delta b \leq O(\lambda |p|^2)
\]
\[
\delta b \leq o(\lambda |p|^2),
\]

as $|p| \to \infty$, uniformly for $x \in \Omega$ and bounded $z$. Then we have, by Theorems 10.3, 13.8, 14.1, 15.2 and 15.5.

\textbf{Theorem 15.10.} \emph{Let $\Omega$ be a bounded domain in $\mathbb{R}^n$ and suppose that the operator $Q$ is elliptic in $\overline{\Omega}$ with coefficients $a^{ij}$, $b \in C^1(\overline{\Omega} \times \mathbb{R} \times \mathbb{R}^n)$ satisfying the structure conditions (15.81) or (15.53) together with condition (10.10) (or (10.36)). Then, if $\partial\Omega \in C^{2,\gamma}$ and $\varphi \in C^{2,\gamma}(\overline{\Omega})$, $0 < \gamma < 1$, there exists a solution $u \in C^{2,\gamma}(\overline{\Omega})$ of the Dirichlet problem $Qu = 0$ in $\Omega$, $u = \varphi$ on $\partial\Omega$.}

In general, if condition (10.10) (or (10.36)) is not assumed in the hypotheses of Theorem 15.10, then the Dirichlet problem $Qu = 0$ in $\Omega$, $u = \varphi$ on $\partial\Omega$, is solvable

% PDF page 393
\source{gilbarg-trudinger:page-0393}

\providecommand{\R}{\mathbb{R}}
\providecommand{\diver}{\operatorname{div}}

provided the family of solutions of the problems (13.42) is uniformly bounded in
\(\Omega\). Note that by Theorem 10.1 the solution is unique if the coefficients \(a^{ij}\) are
independent of \(z\) and the coefficient \(b\) satisfies \(D_z b \leq 0\). In this case, the conditions
\(\delta a^{ij}=o(\lambda), \delta b=o(\lambda |p|^2)\) in (15.81) would be implied by the conditions \(D_x a^{ij}=O(\lambda)\),
\(D_x b=O(\lambda |p|^2)\).

\medskip

\noindent\((\mathrm{ii})\) \emph{Uniformly-elliptic equations in the divergence form} (15.56). Here we assume
the structure equations

\begin{equation}\tag{15.82}
\begin{aligned}
|p|^\tau &\leq O(\lambda)\\
D_p A &= O(|p|^\tau),\\
|p|\,D_z A,\ D_x A,\ B &= O(|p|^{\tau+2})
\end{aligned}
\end{equation}

as \(|p| \to \infty\), uniformly for \(x \in \Omega\) and bounded \(z\), with \(\tau > -1\). We then have, by
Theorems 10.9, 13.8, 14.1 and 15.9.

\medskip

\noindent\textbf{Theorem 15.11.} \emph{Let \(\Omega\) be a bounded domain in \(\R^n\) and suppose that the operator \(Q\) is
elliptic in \(\overline{\Omega}\) with coefficients \(A^i \in C^{1,\gamma}(\overline{\Omega}\times \R \times \R^n)\), \(i=1,\ldots,n\), \(B \in C^\gamma(\overline{\Omega}\times \R \times \R^n)\),
\(0<\gamma<1\), satisfying the structure conditions (15.82) together with the hypotheses
of Theorem 10.9 for \(\alpha=\tau+2\). Then, if \(\partial\Omega \in C^{2,\gamma}\) and \(\varphi \in C^{2,\gamma}(\overline{\Omega})\), there exists a
solution \(u \in C^{2,\gamma}(\overline{\Omega})\) of the Dirichlet problem \(Qu=0\) in \(\Omega\), \(u=\varphi\) on \(\partial\Omega\).}

\medskip

In order to conclude Theorem 15.11 from the estimates in Theorems 10.9, 14.1
and 15.9, we take in Theorem 13.8 the family of Dirichlet problems given by

\begin{equation}\tag{15.83}
Q_\sigma u=\diver\{\sigma A+(1-\sigma)(1+|Du|^2)^{\tau/2}Du\}+\sigma B=0,
\end{equation}

\[
u=\sigma\varphi \text{ on } \partial\Omega,\qquad 0\leq \sigma \leq 1.
\]

Note that the family given by (13.42) does not necessarily satisfy the hypotheses of
Theorem 10.9. As in the previous case, if we do not assume the hypotheses of a
specific maximum principle such as Theorem 10.9, then the Dirichlet problem
\(Qu=0\) in \(\Omega\), \(u=\varphi\) on \(\partial\Omega\) is solvable provided the solutions of a related family of
problems such as (15.83) are uniformly bounded on \(\Omega\). Note that, by Theorem 10.7,
the solution of the problem \(Qu=0\), \(u=\varphi\) on \(\partial\Omega\), is unique provided either the
coefficients \(A\) are independent of \(z\) or the coefficient \(B\) is independent of \(p\) and also
provided \(B\) is non-increasing in \(z\).
    Equations of the form

\begin{equation}\tag{15.84}
Qu=a^{ij}(x,u)D_{ij}u+b(x,u,Du)
\end{equation}

with \(a^{ij}\in C^1(\Omega\times \R)\) can be written in the divergence form (15.57) with

\[
A^i(x,z,p)=a^{ij}(x,z)p_j,
\]
\[
B(x,z,p)=b(x,z,p)-D_z a^{ij}(x,z)p_ip_j-D_{x_i}a^{ij}(x,z)p_j.
\]

% PDF page 394
\source{gilbarg-trudinger:page-0394}

Therefore, using Theorems 10.3, 13.8, 14.1 and 15.9, we have

\medskip

\noindent\textbf{Theorem 15.12.} \emph{Let $\Omega$ be a bounded domain in $\R^n$ and suppose that the operator $Q$
given by (15.84) is elliptic in $\overline{\Omega}$ with coefficients
\[
 a^{ij} \in C^1(\overline{\Omega} \times \R), \qquad b \in C^\gamma(\overline{\Omega} \times \R \times \R^n), \quad 0 < \gamma < 1,
\]
satisfying $b = O(|p|^2)$ as $|p| \to \infty$, uniformly for $x \in \Omega$ and bounded $z$, together with
condition (10.10) (or 10.36). Then, if $\partial \Omega \in C^{2,\gamma}$ and $\varphi \in C^{2,\gamma}(\overline{\Omega})$, there exists a
solution $u \in C^{2,\gamma}(\overline{\Omega})$ of the Dirichlet problem $Qu = 0$ in $\Omega$, $u = \varphi$ on $\partial\Omega$.}

\medskip

\noindent\textbf{(iii)} \emph{Non-uniformly elliptic equations in general domains.} Let us now assume that
the coefficients of equation (15.1) can be decomposed according to (15.23) such that
the following structure conditions are fulfilled

\begin{equation}\tag{15.85}
\begin{aligned}
|p|\, a^{ij},\, b &= O(\mathcal{E}), \\
\delta a^{ij}_* &= O\!\left(\sqrt{\lambda_*\,\mathcal{E}/|p|}\right), \\
\delta a_*^{ij} &= o\!\left(\sqrt{\lambda_*\,\mathcal{E}/|p|}\right), \\
\delta b &\leq O(\mathcal{E}), \\
\delta b &\leq o(\mathcal{E}),
\end{aligned}
\end{equation}

as $|p| \to \infty$, uniformly for $x \in \Omega$ and bounded $z$. Then, by combining Theorems
10.3, 13.8, 14.1 and 15.2, we have

\medskip

\noindent\textbf{Theorem 15.13.} \emph{Let $\Omega$ be a bounded domain in $\R^n$ with coefficients $a^{ij}$, $b \in$
$C^1(\overline{\Omega} \times \R \times \R^n)$ satisfying the structure conditions (15.85) together with condition
(10.10) (or (10.36)). Then, if $\partial\Omega \in C^{2,\gamma}$ and $\varphi \in C^{2,\gamma}(\overline{\Omega})$, $0 < \gamma < 1$, there exists a
solution $u \in C^{2,\gamma}(\overline{\Omega})$ of the Dirichlet problem $Qu = 0$ in $\Omega$, $u = \varphi$ on $\partial\Omega$.}

\medskip

Theorem 15.13 is clearly an extension of Theorem 15.10 and the remarks
following the latter are also pertinent here. Note that when a decomposition of the
form (15.2) is valid we have $\delta a_*^{ij} = 0$, and moreover if $a_*^{ij}(p) = a_*^{ij}(p/|p|)$ we also have
$\delta a_*^{ij} = 0$.

\medskip

\noindent\textbf{(iv)} \emph{Non-uniformly elliptic equations in convex domains.} Let us now assume
that the decomposition (15.2) is valid and

\begin{equation}\tag{15.86}
\begin{aligned}
b &= o(|p|\Lambda), \\
\delta b &\leq O\!\left(b^2/|p|^2\,\mathcal{T}_*\right), \qquad (\mathcal{T}_* = \operatorname{trace}[a_*^{ij}]),
\end{aligned}
\end{equation}

as $|p| \to \infty$, uniformly in $x \in \Omega$ and bounded $z$. Then from Theorems 13.8, 15.1
and Corollary 14.5 we have

\medskip

\noindent\textbf{Theorem 15.14.} \emph{Let $\Omega$ be a uniformly convex domain in $\R^n$ and suppose that the
operator $Q$ is elliptic in $\Omega$ with coefficients $a^{ij}$, $b \in C^1(\overline{\Omega} \times \R \times \R^n)$ satisfying the}

% PDF page 395
\source{gilbarg-trudinger:page-0395}

\textbf{15.5. Selected Existence Theorems} \hfill 383

\smallskip

\emph{structure conditions (15.86). Then, if $\partial \Omega \in C^{2,\gamma}$ and $\varphi \in C^{2,\gamma}(\overline{\Omega})$, $0<\gamma<1$, there exists a solution $u \in C^{2,\gamma}(\overline{\Omega})$ of the Dirichlet problem $Qu=0$ in $\Omega$, $u=\varphi$ on $\partial \Omega$, provided the family of $C^{2,\gamma}(\overline{\Omega})$ solutions of the problems}

\[
\tag{15.87}
Q_\sigma u = a^{ij}(x,u,Du)D_{ij}u + \sigma b(\sigma x + (1-\sigma)x_0,\sigma u,Du)=0,
\]
\[
u=\sigma\varphi \text{ on } \partial \Omega,\qquad 0\le \sigma \le 1,
\]

\noindent \emph{is uniformly bounded in $\Omega$, for some fixed $x_0\in\Omega$.}

\medskip

An existence theorem for non-uniformly elliptic equations in the divergence form (15.57) follows in a similar fashion when Theorem 15.6 is used in place of Theorem 15.1 above. We assume that the coefficients $A^i$ and $B$ in (15.56) satisfy for some $\tau\in\mathbb{R}$, the conditions

\[
|p|^\tau \le O(\lambda),
\]
\[
\tag{15.88}
|p|\,D_z A,\, D_x A,\, B = o(|p|^{\tau+1})
\]

\noindent as $|p|\to\infty$, uniformly for $x\in\Omega$ and bounded $z$. Then we have

\medskip

\textbf{Theorem 15.15.} \emph{Let $\Omega$ be a uniformly convex domain in $\mathbb{R}^n$ and suppose that the operator $Q$ is elliptic in $\overline{\Omega}$ with coefficients $A^i\in C^1(\overline{\Omega}\times\mathbb{R}\times\mathbb{R}^n)$, $i=1,\ldots,n$, $B\in C^1(\overline{\Omega}\times\mathbb{R}\times\mathbb{R}^n)$, $0<\gamma<1$, satisfying the structure conditions (15.88). Then, if $\partial\Omega\in C^{2,\gamma}$ and $\varphi\in C^{2,\gamma}(\overline{\Omega})$, there exists a solution $u\in C^{2,\gamma}(\overline{\Omega})$ of the Dirichlet problem $Qu=0$ in $\Omega$, $u=\varphi$ on $\partial\Omega$, provided the family of $C^{2,\gamma}(\overline{\Omega})$ solutions of the problems (13.42) is uniformly bounded in $\Omega$.}

\medskip

\textbf{(v)} Problems with boundary curvature conditions. Let us now consider operators that are decomposed according to both (14.43) and (15.23). In particular, we shall assume that

\[
a^{ij}(x,z,p)=\Lambda a_\infty^{ij}(x,p/|p|)+a_0^{ij}(x,z,p)
\]
\[
\tag{15.89}
= a_*^{ij}(x,z,p)+\frac12\bigl[p_ic_j(x,z,p)+c_i(x,z,p)p_j\bigr],
\]
\[
b(x,z,p)=|p|\Lambda b_\infty(x,z,p/|p|)+b_0(x,z,p)
\]

\noindent where $a_\infty^{ij}\in C^1(\overline{\Omega}\times B(0))$, $a_*^{ij}, b_\infty\in C^1(\overline{\Omega}\times\mathbb{R}\times B(0))$, $i,j=1,\ldots,n$, the matrices $[a_\infty^{ij}]$, $[a_*^{ij}]$ are non-negative and symmetric, and $b_\infty$ is non-increasing with respect to $z$. We shall impose the following structure conditions

\[
\delta\mathcal{E}\le \mathcal{E}+o(\mathcal{E}),
\]
\[
\delta\mathcal{E},\,\delta b,\;(\delta-1)b_0\le O(\mathcal{E}),
\]
\[
\tag{15.90}
\delta a_*^{ij}=O\!\left(\frac{\sqrt{\mathcal{E}}}{|p|}\right),
\]
\[
a_0^{ij}=o(\Lambda)
\]
\[
b_0=o(|p|)
\]

% PDF page 396
\source{gilbarg-trudinger:page-0396}


as $|p|\to\infty$, uniformly for $x\in\Omega$ and bounded $z$. Then we have, by Theorems 13.8,
14.9 and 15.2,

\medskip

\textbf{Theorem 15.16.} \emph{Let $\Omega$ be a bounded domain in $\mathbb{R}^n$ and suppose that the operator $Q$
is elliptic in $\overline{\Omega}$ with coefficients $a^{ij}$, $b\in C^1(\overline{\Omega}\times\mathbb{R}\times\mathbb{R}^n)$ satisfying the structure
conditions (15.89), (15.90). Then, if $\partial\Omega\in C^{2,\gamma}$, $\varphi\in C^{2,\gamma}(\overline{\Omega})$, $0<\gamma<1$, and the
inequalities}

\[
\tag{15.91}
\mathcal{K}^{+}>b_\infty(y,\varphi(y),v), \qquad \mathcal{K}^{-}>-b_\infty(y,\varphi(y),-v)
\]

\emph{hold at each point $y\in\partial\Omega$, where $v$ is the unit inner normal at $y$ and $\mathcal{K}^{+}$, $\mathcal{K}^{-}$, given by
(14.44), are non-negative, it follows that the Dirichlet problem $Qu=0$ in $\Omega$, $u=\varphi$
on $\partial\Omega$, is solvable provided the family of $C^{2,\gamma}(\overline{\Omega})$ solutions of the problems (13.42) is
uniformly bounded in $\Omega$.}

\medskip

In order to permit the non-strict inequalities in (15.91) we need to strengthen
the structure conditions (15.90) so that

\[
\delta\mathcal{E}\le \mathcal{E}+o(\mathcal{E}),
\]
\[
% [illegible] leading symbol before \delta\mathcal{E} in the printed scan
\delta\mathcal{E},\,\delta b,\,\bar{\delta}b_0\le O(\mathcal{E}),
\]
\[
\tag{15.92}
\delta a_*^{ij}=O\!\left(\sqrt{\mathcal{E}}/|p|\right),
\]
\[
a_0^{ij}=O(\mathcal{E}/|p|),
\]
\[
b_0=O(\mathcal{E}),
\]

\noindent as $|p|\to\infty$, uniformly for $x\in\Omega$ and bounded $z$. Then we have, by Theorems 13.8,
14.6, 15.2,

\medskip

\textbf{Theorem 15.17.} \emph{Let $\Omega$ be a bounded domain in $\mathbb{R}^n$ and suppose that the operator $Q$
is elliptic in $\overline{\Omega}$ with coefficients $a^{ij}$, $b\in C^1(\overline{\Omega}\times\mathbb{R}\times\mathbb{R}^n)$ satisfying the structure
conditions (15.89), (15.92). Then, if $\partial\Omega\in C^{2,\gamma}$, $\varphi\in C^{2,\gamma}(\overline{\Omega})$, $0<\gamma<1$, and the
inequalities}

\[
\tag{15.93}
\mathcal{K}^{+}\ge b_\infty(y,\varphi(y),v), \qquad \mathcal{K}^{-}\ge -b_\infty(y,\varphi(y),-v), \qquad \mathcal{K}^{\pm}\ge 0
\]

\emph{hold at each point $y\in\partial\Omega$, it follows that the Dirichlet problem $Qu=0$ in $\Omega$, $u=\varphi$ on
$\partial\Omega$, is solvable provided the family of $C^{2,\gamma}(\overline{\Omega})$ solutions of the problems (13.42) is
uniformly bounded in $\Omega$.}

\bigskip

\textbf{15.6. Existence Theorems for Continuous Boundary Values}

\medskip

By means of the interior estimates, Theorems 15.3 and 15.6, certain of the existence
theorems of the preceding section can be extended to hold when the function $\varphi$ is
only assumed to be continuous on $\partial\Omega$. The basic procedure is to approximate the


% PDF page 397
\source{gilbarg-trudinger:page-0397}

function $\varphi$ uniformly on $\partial\Omega$ by functions $\varphi_m \in C^{2,\gamma}(\Omega)$ and solve the Dirichlet problems, $Qu_m = 0$ in $\Omega$, $u = \varphi_m$ on $\partial\Omega$, for functions $u_m \in C^{2,\gamma}(\overline{\Omega})$. The interior estimates, Theorems 15.3 and 15.8, on combination with the interior Hölder derivative estimates, Theorems 13.1 or 13.3 and the Schauder interior estimate, Theorem 6.2, guarantee that a subsequence of the sequence $\{u_m\}$ converges uniformly on compact subsets of $\Omega$, together with its first and second derivatives, to a function $u \in C^{2,\gamma}(\Omega)$ satisfying $Qu=0$ in $\Omega$. The modulus of continuity estimates, Theorem 14.15, ensure that $u \in C^0(\overline{\Omega}) \cap C^{2}(\Omega)$ and moreover that $u=\varphi$ on $\partial\Omega$. For this procedure to work we also require that the sequence $\{u_m\}$ be uniformly bounded, that is we require a maximum principle as in Chapter 10. By approximation of the domain $\Omega$ by $C^{2,\gamma}$ domains we can also, through Theorem 14.15, remove the restriction that $\partial\Omega \in C^{2,\gamma}$. Let us now state two existence theorems for general domains which can be obtained through this procedure. In Section 16.3, we shall consider analogous results for the minimal surface and prescribed mean curvature equations.

\medskip

\textbf{Theorem 15.18.} \emph{Let $\Omega$ be a bounded domain in $\R^n$ satisfying an exterior sphere condition at each point of the boundary $\partial\Omega$. Let $Q$ be an elliptic operator on $\Omega$ with coefficients $a^{ij}, b \in C^{1}(\Omega \times \R \times \R^n)$ satisfying the hypotheses of Theorems 15.3 (or 15.5) and 14.1 together with condition (10.10) (or (10.36)). Then for any function $\varphi \in C^0(\partial\Omega)$, there exists a solution $u \in C^0(\overline{\Omega}) \cap C^{2}(\Omega)$ of the Dirichlet problem $Qu = 0$ in $\Omega$, $u = \varphi$ on $\partial\Omega$.}

\medskip

For equations in the divergence form (15.57), we have

\medskip

\textbf{Theorem 15.19.} \emph{Let $\Omega$ be a bounded domain in $\R^n$ satisfying an exterior sphere condition at each point of the boundary $\partial\Omega$. Let $Q$ be a divergence structure operator with coefficients $A^i \in C^{1,\gamma}(\Omega \times \R \times \R^n)$, $i=1,\dots,n$, $B \in C^{\gamma}(\Omega \times \R \times \R^n)$, $0 < \gamma < 1$, satisfying the hypotheses of Theorem 15.8, together with the hypotheses of Theorem 10.9 for $\alpha = \gamma + 2$. Then, for any function $\varphi \in C^0(\partial\Omega)$, there exists a solution $u \in C^0(\overline{\Omega}) \cap C^{2,\gamma}(\Omega)$ of the Dirichlet problem $Qu=0$ in $\Omega$, $u=\varphi$ on $\partial\Omega$.}

\bigskip

\textbf{Notes}

The essential ideas in the maximum principle approach to gradient estimates, presented in Sections 15.1 and 15.2, go back to Bernstein [BE 3, 6]. Bernstein's method was substantially developed by Ladyzhenskaya [LA] and Ladyzhenskaya and Ural'tseva [LU 2, 4] to yield both global and interior gradient estimates for uniformly elliptic equations. Later Serrin [SE 3], by exploiting representations of the form (15.2), extended these results to encompass equations with similar features to the prescribed mean curvature equation (10.7). Theorems 15.2 and 15.3 are very close to results in [LU 5, 6] (see also [IV 1, 2]), although they, as well as Theorem 15.1, are formulated similarly to the statements in [SE 4]. Rather than use multipliers $r, s, t$, the authors in [LU 5, 6] point out that their hypotheses need only be satisfied by an operator equivalent to the given one. Our treatment differs from

% PDF page 398
\source{gilbarg-trudinger:page-0398}

those in [LU 5, 6] and [SE 4], in that we consider $C^2(\Omega)$ solutions instead of $C^3(\Omega)$
solutions. Also, our proofs and results will continue to be valid for solutions in the
space $C^{0,1}(\overline{\Omega}) \cap W^{2,2}(\Omega)$.

The global gradient bound for solutions of divergence structure equations,
Theorem 15.6, is due to Trudinger [TR 3] while the estimates of Theorems 15.7,
15.8 and 15.9 for uniformly elliptic divergence structure equations are due to
Ladyzhenskaya and Ural'tseva [LU 1, 2]. Our proof of Theorem 15.8 differs in
some aspects from that in [LU 2]. For further gradient estimates for non-
uniformly elliptic, divergence structure equations, see [IO], [OS 1, 2]. Some of the
existence theorems in Sections 15.5 and 15.6 are already formulated in the literature,
see for example [LU 4], [SE 3]. The survey paper [ED] provides a clear account
of various aspects of the execution of the existence procedure. We note also that
many of the above existence theorems extend to $C^{1,\alpha}$ boundary data, $0<\alpha<1$;
(see [LB 1]).

Finally, we remark that Lemma 15.4 resolved positively the long standing
problem as to whether the natural conditions (that is, the case $\theta = 1, r=s=t_i=0$
in (15.53)) alone suffice for interior and global gradient bounds; (see [LU 7],
[TR 12]).

\medskip

\textbf{Problems}

\medskip

\textbf{15.1.} Show that Theorem 15.2 continues to hold for positive $a$ and $c$ provided
$\operatorname{osc}_{\Omega} u$ is sufficiently small.

\medskip

\textbf{15.2.} Show that Theorems 15.2 and 15.3 continue to hold for positive $a$ and $c$
provided $b<-2\sqrt{ac}$.

\medskip

\textbf{15.3.} Show that Theorem 15.6 continues to hold if the structure condition (15.60)
is replaced by

\[
(15.94)\qquad |p|\,|D_z A, D_x A, B| = o(|p|^\gamma)\qquad \text{as } |p| \to \infty,
\]

where $\gamma = \tau + 1 + (\tau + 2)/n$, provided $\tau > -1$ and $u=0$ on $\partial\Omega$.

\medskip

\textbf{15.4.} Show that Theorem 15.6 continues to hold if the structure condition (15.60)
is replaced by

\[
(15.95)\qquad |p|\,|D_z A, D_x A, B| = o(|p|^\gamma)\qquad \text{as } |p| \to \infty,
\]

where $\gamma = \tau + 1 + 1/n$, provided $\tau \le -1$ and $u=0$ on $\partial\Omega$.

% PDF page 399
\source{gilbarg-trudinger:page-0399}

\begin{quote}
\textbf{Problems}
\end{quote}

\textbf{15.5.} Show that an interior gradient bound is valid for $C^2(\Omega)$ solutions of equation
\[
(15.56)\quad \text{provided that the structure conditions } (15.60),
\]
\[
(15.64)'\qquad D_p A = O(|p|^{\tau+\sigma}),
\]
\[
(15.66)'\qquad g = o(|p|^{\tau+2})
\]
hold for $\tau > -1$, $\sigma < 1$.

\textbf{15.6.} Consider the Dirichlet problem
\[
(15.96)\qquad \varepsilon \Delta u + g(Du) = f(x) \quad \text{in } \Omega,\qquad u = 0 \quad \text{on } \partial\Omega,
\]
in a bounded convex domain $\Omega \subset \mathbb{R}^n$, where $\varepsilon > 0$, $g \in C^1(\mathbb{R}^n)$, $f \in C^1(\overline{\Omega})$, $\sqrt{|p|}/g = O(1)$ as $|p| \to \infty$, and $g(0) \leq f(x)$ for all $x \in \Omega$. Using Theorem 15.1, together with linear barrier functions, deduce the unique solvability of (15.96), for sufficiently small $\varepsilon$, in the space $C^{0,1}(\overline{\Omega}) \cap C^2(\Omega)$. By letting $\varepsilon \to 0$, show that there exists a $C^{0,1}(\overline{\Omega})$ solution of the \textit{first order} Dirichlet problem
\[
(15.97)\qquad g(Du) = f(x) \quad \text{a.e. } (\Omega),\qquad u = 0 \quad \text{on } \partial\Omega,
\]
and using suitable approximation show that the above conditions on $f$ and $g$ can be replaced by $g \in C^0(\mathbb{R}^n)$, $g \to \infty$ as $|p| \to \infty$, $f \in L^\infty(\Omega)$, $g(0) \leq f(x)$ for all $x \in \Omega$. Find a similar result for arbitrary domains satisfying exterior sphere conditions. In a certain sense the solution of (15.97) thus obtained is unique. (Cf. [LP 3, 5].)

% PDF page 400
\source{gilbarg-trudinger:page-0400}


Chapter 16

Equations of Mean Curvature Type

In this chapter we focus attention on both the prescribed mean curvature equation,

(16.1) \qquad $\mathfrak{M}u=(1+|Du|^2)\Delta u-D_i uD_j uD_{ij}u=nH(1+|Du|^2)^{3/2},$

and a related family of equations in two variables. Our main concern is with
interior derivative estimates for solutions. We shall see that not only can interior
gradient bounds be established for solutions of these equations but that also their
non-linearity leads to strong second derivative estimates which distinguish them
from uniformly elliptic equations such as Laplace's equation. In particular we shall
derive an extension of the classical result of Bernstein that a $C^2(\R^2)$ solution of the
minimal surface equation in $\R^2$ must be a linear function (Theorem 16.12).

The approach to interior gradient bounds, in this chapter, differs considerably
from those in Chapter 15 (although it does have some features in common with the
divergence structure case of Section 15.4). An interior gradient bound for a solution
of equation (16.1), Theorem 16.5, is derived through a consideration of the tan-
gential gradient and Laplace operators on the hypersurface in $\R^{n+1}$ given by the
graph of the solution $u$. The basic estimates on hypersurfaces are supplied in
Section 16.1.

The study of the general class of equations of prescribed mean curvature type in
two variables is taken up in Section 16.4. Interior first and second derivative
estimates (Theorems 16.20 and 16.21) arise through a treatment in Section 16.5 of
quasiconformal mappings between surfaces in $\R^3$, which extends that of Section
12.1.

Sections 16.4 to 16.8 and a portion of Section 16.1 of this chapter were written
in collaboration with L. M. Simon and they are essentially his contribution.

\medskip

16.1. Hypersurfaces in $\R^{n+1}$

 A subset $\mathcal G$ of $\R^{n+1}$ is called a $C^k$ hypersurface in $\R^{n+1}$ if locally $\mathcal G$ can be represented
as the graph of a $C^k$ function over an open subset of $\R^n$. Our concern here is with
$C^2$ hypersurfaces and for convenience we shall assume that $\mathcal G$ can be globally
represented as the level surface of a $C^2$ function, that is, there exists an open subset
